\chapter{ПРОЕКТИРОВАНИЕ И РЕАЛИЗАЦИЯ СИСТЕМЫ}

\section{Техническое задание}

\subsection{Назначение разработки}

Система BioSonification предназначена для генерации символической музыки на основе биологических последовательностей. Система должна преобразовывать FASTA-файлы в двухдорожечные MIDI-композиции, используя биологические признаки как управляющий сигнал.

\subsection{Функциональные требования}

\begin{enumerate}
\item Чтение FASTA-файлов с DNA, RNA и protein последовательностями
\item Фрагментация длинных последовательностей
\item Извлечение биологических признаков (состав, энтропия, белковые свойства)
\item Загрузка и структурирование MIDI-корпуса
\item Обучение модели гармонии
\item Обучение модели мелодии
\item Генерация двухдорожечного MIDI
\item Сохранение метаданных генерации
\end{enumerate}

\subsection{Нефункциональные требования}

\begin{itemize}
\item Обучение на GPU NVIDIA RTX 2060 6GB
\item Поддержка mixed precision для экономии памяти
\item Воспроизводимость результатов (фиксированный seed)
\item Сохранение checkpoint для повторного использования
\end{itemize}

\section{Архитектура системы}

Система построена по модульному принципу и включает следующие компоненты:

\begin{itemize}
\item \textbf{Биологический энкодер} — извлечение признаков из FASTA
\item \textbf{Музыкальный процессор} — структурирование MIDI-корпуса
\item \textbf{Pairing модуль} — сопоставление биологических и музыкальных сегментов
\item \textbf{Harmony model} — генерация аккордовой сетки
\item \textbf{Melody model} — генерация мелодии поверх гармонии
\item \textbf{MIDI renderer} — создание выходного файла
\end{itemize}

% Общая схема пайплайна представлена на рисунке~\ref{fig:architecture}.

% \begin{figure}[htbp]
% \centering
% \includegraphics[width=0.9\textwidth]{pipeline_architecture.png}
% \caption{Архитектура системы BioSonification}
% \label{fig:architecture}
% \end{figure}

\section{Биологический энкодер}

Модуль \texttt{bio\_music\_pipeline/v2/bio.py} реализует извлечение признаков из биологических последовательностей.

\subsection{Поддерживаемые типы последовательностей}

Система автоматически определяет тип последовательности:
\begin{itemize}
\item DNA — содержит A, C, G, T
\item RNA — содержит A, C, G, U
\item Protein — содержит аминокислоты
\end{itemize}

\subsection{Фрагментация}

Длинные последовательности разбиваются на фрагменты с параметрами:
\begin{itemize}
\item \texttt{fragment\_length} = 1800 bp
\item \texttt{fragment\_stride} = 900 bp
\item \texttt{max\_fragments\_per\_record} = 24
\end{itemize}

\subsection{Извлекаемые признаки}

\textbf{Нуклеотидные признаки:}
\begin{itemize}
\item Частоты оснований (A, C, G, T/U)
\item Энтропия последовательности
\item GC-content
\item GC/AT skew
\item Кодоновая энтропия
\item k-mer статистика (k=1,2,3)
\end{itemize}

\textbf{Белковые признаки:}
\begin{itemize}
\item Перевод через longest ORF
\item Ароматичность (Biopython ProtParam)
\item Индекс нестабильности
\item Изоэлектрическая точка
\item GRAVY (гидрофобность)
\item Молекулярная масса
\item Состав аминокислот
\end{itemize}

\subsection{Выходные данные}

Энкодер возвращает структуру \texttt{BioEncodingResult}:
\begin{itemize}
\item \texttt{vector} — 256-мерный embedding
\item \texttt{control\_profile} — 6-мерный профиль (tempo, density, harmony change, register, complexity, mode)
\item \texttt{tonic\_pc\_hint} — подсказка тонального центра
\end{itemize}

\section{Музыкальное представление}

Модуль \texttt{structured\_music.py} реализует структурированное представление музыки.

\subsection{Структуры данных}

\textbf{HarmonyBar} — аккорд в такте:
\begin{itemize}
\item \texttt{bar\_index} — номер такта
\item \texttt{root\_pc} — основной тон (0-11)
\item \texttt{quality} — качество аккорда (maj, min, dim, aug, dom7, maj7, min7, sus2, sus4, power)
\item \texttt{hold} — продолжение предыдущего аккорда
\end{itemize}

\textbf{MelodyEvent} — событие мелодии:
\begin{itemize}
\item \texttt{bar\_index} — номер такта
\item \texttt{position\_step} — позиция в такте
\item \texttt{duration\_steps} — длительность
\item \texttt{relative\_pc} — высота относительно аккорда
\item \texttt{octave} — октава
\end{itemize}

\subsection{Токенизация}

Система использует два токенизатора:
\begin{itemize}
\item \texttt{HarmonyTokenizer} — кодирует последовательность аккордов
\item \texttt{MelodyTokenizer} — кодирует мелодические события с учётом гармонии
\end{itemize}

\section{Pairing биологических и музыкальных сегментов}

Модуль \texttt{structured\_pairing.py} сопоставляет биологические фрагменты с музыкальными сегментами через пространство дескрипторов.

\subsection{Музыкальные дескрипторы}

Для каждого музыкального сегмента вычисляются:
\begin{itemize}
\item Темп (BPM)
\item Плотность нот на такт
\item Частота смены аккордов
\item Средний регистр
\item Сложность гармонии
\item Модальность (мажор/минор)
\end{itemize}

\subsection{Калибровка}

Биологический \texttt{control\_profile} калибруется под распределение музыкального корпуса, чтобы значения попадали в реалистичный диапазон.

\subsection{Top-k matching}

Для каждого биологического фрагмента подбираются k=5 ближайших музыкальных сегментов по евклидову расстоянию в пространстве дескрипторов.

\section{Нейросетевая модель}

Модуль \texttt{structured\_model.py} реализует \texttt{BioConditionedSequenceModel} — autoregressive Transformer с bio conditioning.

\subsection{Архитектура}

Параметры актуальной модели, используемой веб-интерфейсом:
\begin{itemize}
\item \texttt{d\_model} = 384
\item \texttt{n\_layers} = 6
\item \texttt{n\_heads} = 6
\item \texttt{dim\_feedforward} = 1536
\item \texttt{dropout} = 0.20
\item \texttt{bio\_memory\_tokens} = 4
\end{itemize}

\subsection{Bio conditioning}

Биологический вектор проецируется в пространство модели через MLP и представляется как memory tokens, которые участвуют в cross-attention.

\subsection{Использование модели}

Одна и та же архитектура используется дважды:
\begin{enumerate}
\item Как harmony model — генерирует последовательность аккордов
\item Как melody model — генерирует мелодию с harmony prefix
\end{enumerate}

\section{Процесс обучения}

\subsection{Конфигурация}

Актуальная модель веб-интерфейса использует конфиг \texttt{pipeline\_v2\_medium\_rtx2060\_fast.json}:
\begin{itemize}
\item Batch size: 8
\item Gradient accumulation: 4 (effective batch size 32)
\item Learning rate: 0.0003
\item Weight decay: 0.05
\item Mixed precision: enabled
\item Optimizer: AdamW
\item Segment length: 4 такта
\end{itemize}

\subsection{Данные}

\begin{itemize}
\item Bio sequences/fragments: 3746
\item Music segments: 25940 (POP909)
\item Train pairs: 15925
\item Validation pairs: 1870
\item Test pairs: 935
\end{itemize}

\subsection{Обучение harmony model}

\begin{itemize}
\item Epochs: 10
\item Best validation loss: 0.1454 (epoch 2)
\item Test loss: 0.1389
\item Training time: около 21 минуты
\end{itemize}

\subsection{Обучение melody model}

\begin{itemize}
\item Epochs: 15
\item Best validation loss: 0.1566 (epoch 7)
\item Test loss: 0.1650
\item Training time: около 31 минуты
\end{itemize}

\section{Реализация}

Проект реализован на Python 3.10+ с использованием:
\begin{itemize}
\item PyTorch 2.0+ — нейросетевой фреймворк
\item Biopython — биоинформатические инструменты
\item music21 — музыкальный анализ
\item mido — работа с MIDI
\item Flask — веб-интерфейс
\end{itemize}

Основные модули:
\begin{itemize}
\item \texttt{train\_bio\_music\_v2.py} — CLI обучения
\item \texttt{generate\_from\_fasta\_v2\_fragmented.py} — CLI генерации
\item \texttt{bio\_music\_pipeline/v2/} — основная логика
\item \texttt{web/app.py} — веб-интерфейс
\end{itemize}

\section*{Выводы по главе 2}

Во второй главе описано проектирование и реализация системы BioSonification. Разработана иерархическая архитектура Bio→Harmony→Melody. Реализованы модули извлечения биологических признаков, структурирования музыки, pairing и нейросетевой генерации.

Система обучена на корпусе POP909 и геномных данных RefSeq. Актуальная модель веб-интерфейса показала validation loss 0.1454 для гармонии и 0.1566 для мелодии. Обучение выполнено на GPU RTX 2060 6GB с использованием mixed precision.
